\documentclass[a4paper,11pt]{article}
\usepackage{lmodern}
\usepackage[T1]{fontenc}
\usepackage[utf8]{inputenc}
\usepackage{amsmath,amssymb,amsthm}
\usepackage{longtable,booktabs}
\usepackage{hyperref}
\usepackage[margin=25mm]{geometry}

\newtheorem{theorem}{Theorem}[section]
\newtheorem{proposition}[theorem]{Proposition}
\newtheorem{lemma}[theorem]{Lemma}
\newtheorem{corollary}[theorem]{Corollary}
\newtheorem{definition}[theorem]{Definition}
\newtheorem{remark}[theorem]{Remark}
\newtheorem{observation}[theorem]{Observation}

\providecommand{\tightlist}{\setlength{\itemsep}{0pt}\setlength{\parskip}{0pt}}

\begin{document}

\section{Formulation of Structural Requirements for a Theory of
Everything}\label{formulation-of-structural-requirements-for-a-theory-of-everything}

\subsection{--- Structural Constraints Based on Definitions of
Information and Operations
---}\label{structural-constraints-based-on-definitions-of-information-and-operations}

\textbf{Author:} Noriaki Kihara (WF System Co., Ltd.)

\textbf{Date:} April 2026

\textbf{DOI:}
\href{https://doi.org/10.5281/zenodo.19601592}{10.5281/zenodo.19601592}

\begin{center}\rule{0.5\linewidth}{0.5pt}\end{center}

\subsection{0. Scope of Claims (Must
Read)}\label{scope-of-claims-must-read}

This paper is a thought experiment that formulates the structural
requirements a Theory of Everything must satisfy, based on rigorous
definitions of information and operations.

\subsubsection{What This Paper Claims}\label{what-this-paper-claims}

\begin{enumerate}
\def\labelenumi{\arabic{enumi}.}
\tightlist
\item
  To provide rigorous definitions of information \(I\) and operation
  \(F^*\) as the most fundamental starting point for a Theory of
  Everything.
\item
  To show that causal structure inevitably emerges from the requirement
  of cognitive accessibility.
\item
  To enumerate the structural requirements that a Theory of Everything
  must satisfy, built upon these definitions and the causal structure.
\item
  To organize the consequences derived from these requirements.
\end{enumerate}

\subsubsection{What This Paper Does Not
Claim}\label{what-this-paper-does-not-claim}

\begin{enumerate}
\def\labelenumi{\arabic{enumi}.}
\tightlist
\item
  This paper is not a theory of physics. It makes no claims about the
  actual universe.
\item
  This paper does not presuppose any specific geometric construction
  (such as central projection). The structural requirements are
  formulated as general demands independent of geometry.
\item
  This paper makes no empirical predictions.
\end{enumerate}

\begin{center}\rule{0.5\linewidth}{0.5pt}\end{center}

\subsection{1. Definitions}\label{definitions}

This chapter provides the most fundamental definitions that serve as the
starting point for a Theory of Everything. Only \textbf{information} and
\textbf{operations} are defined here; concepts such as causality, time,
space, and matter are not presupposed at all.

The definitions in this chapter are constructed as the minimal axiom
system that is necessary and sufficient as a foundation for physical
theories, referencing the mathematical foundations of information
{[}1{]} and the framework of computability theory {[}2{]}.

\subsubsection{\texorpdfstring{1.1 Definition of Information
\(I\)}{1.1 Definition of Information I}}\label{definition-of-information-i}

\textbf{Information \(I\)} refers to any information that can be
expressed numerically. Its format is arbitrary: scalar values, complex
values, matrices, tensors, etc. However, information \(I\) must satisfy
the following properties:

\begin{enumerate}
\def\labelenumi{\arabic{enumi}.}
\tightlist
\item
  \textbf{Existence of null state}: A state in which no information
  exists, \(I = \mathrm{null}\), can be defined.
\item
  \textbf{Decidability of null}: For any information \(I\), it is
  decidable whether \(I = \mathrm{null}\) or not. The comparison
  \(\mathrm{null} = \mathrm{null}\) is true.
\item
  \textbf{Definition of type}: Each piece of information belongs to a
  type. For any two pieces of information \(I_a,\, I_b\), it is
  decidable whether they belong to the same type. For example, a real
  number and an array are different types. \(\mathrm{null}\) belongs to
  every type (i.e., for any information \(I\), the type comparison
  between \(\mathrm{null}\) and \(I\) always yields the same type).
\item
  \textbf{Decidability of equivalence}: For any two pieces of
  information \(I_a,\, I_b\) belonging to the same type, it is decidable
  whether \(I_a = I_b\) (whether they are identical information).
\end{enumerate}

\subsubsection{\texorpdfstring{1.2 Definition of Operation
\(F^*\)}{1.2 Definition of Operation F\^{}*}}\label{definition-of-operation-f}

The mapping from information to information

\[
\mathrm{OUT} = F^*(\mathrm{IN})
\]

is defined as an \textbf{operation}. Operation \(F^*\) requires that the
input information and output information belong to the same type.

Operation \(F^*\) has the following internal structure:

\begin{itemize}
\tightlist
\item
  \textbf{Input information \(\mathrm{IN}\)}: Information provided
  externally. It can only be referenced and cannot be modified from
  within \(F^*\).
\item
  \textbf{Internal information \(I\)}: Information held internally by
  \(F^*\). It is observable externally as \(\mathrm{OUT}\) but cannot be
  modified from outside.
\end{itemize}

\paragraph{Initial State}\label{initial-state}

The initial state of operation \(F^*\) is \(I = \mathrm{null}\),
\(\mathrm{OUT} = \mathrm{null}\).

\paragraph{Operation Rules}\label{operation-rules}

Operation \(F^*\) operates according to the following rules:

\begin{enumerate}
\def\labelenumi{\arabic{enumi}.}
\tightlist
\item
  If \(\mathrm{IN} \neq I\), or \(I = \mathrm{null}\):

  \begin{itemize}
  \tightlist
  \item
    Copy input information \(\mathrm{IN}\) to internal information \(I\)
    (\(I \leftarrow \mathrm{IN}\))
  \item
    Output internal information \(I\) to output information
    \(\mathrm{OUT}\)
  \end{itemize}
\item
  If \(\mathrm{IN} = I\):

  \begin{itemize}
  \tightlist
  \item
    Do nothing (\(\mathrm{OUT}\) remains unchanged)
  \end{itemize}
\end{enumerate}

\subsubsection{1.3 Completeness of the
Definitions}\label{completeness-of-the-definitions}

The above definitions (§1.1--§1.2) constitute the minimal definitions
that are necessary and sufficient as a starting point for a Theory of
Everything.

These definitions:

\begin{itemize}
\tightlist
\item
  Do not presuppose causality (no temporal ordering between IN and OUT
  is assumed)
\item
  Do not presuppose time
\item
  Do not presuppose space
\item
  Do not presuppose the number of dimensions
\item
  Do not presuppose any specific physical theory
\item
  Do not presuppose any specific geometry
\item
  Do not restrict the format of information (scalars, vectors, matrices,
  tensors, etc. are all permitted)
\end{itemize}

Under these definitions alone, causality does not exist. A stationary
state of information (a state without change) is described without
contradiction under these definitions.

\begin{center}\rule{0.5\linewidth}{0.5pt}\end{center}

\subsection{2. Cognitive Accessibility and the Emergence of
Causality}\label{cognitive-accessibility-and-the-emergence-of-causality}

\subsubsection{2.1 Requirement of Cognitive
Accessibility}\label{requirement-of-cognitive-accessibility}

Under the definitions of §1 alone, information and operations ``exist''
but there is no subject to distinguish or identify them. The minimal
definition of ``cognition'' is as follows:

\textbf{Cognitive accessibility:} The input information \(\mathrm{IN}\)
and output information \(\mathrm{OUT}\) of operation \(F^*\) can be
distinguished. That is, an \textbf{ordering} exists between
\(\mathrm{IN}\) and \(\mathrm{OUT}\).

\subsubsection{2.2 Inevitable Emergence of
Causality}\label{inevitable-emergence-of-causality}

The moment cognitive accessibility is required, the following inevitably
emerge:

\begin{enumerate}
\def\labelenumi{\arabic{enumi}.}
\tightlist
\item
  \textbf{Emergence of ordering:} A direction arises in which
  \(\mathrm{IN}\) comes first and \(\mathrm{OUT}\) comes after.
\item
  \textbf{Emergence of causality:} The three-stage separation of
  \(\mathrm{IN}\) (cause) → \(F^*\) (process) → \(\mathrm{OUT}\)
  (effect) is established.
\item
  \textbf{Emergence of operation count:} A cumulative count \(n\) (a
  non-negative integer) of how many times operation \(F^*\) has been
  executed becomes definable.
\end{enumerate}

The operation count \(n\) is a \textbf{value local to each instance} of
operation \(F^*\), not a global counter for the entire system. The
absence of a global time parameter is consistent with the ``frozen time
problem'' in general relativity {[}3{]}.

\subsubsection{2.3 Permission of Causality-Free
States}\label{permission-of-causality-free-states}

An important point is that the definitions of §1 also permit states
without causality. That is, when cognitive accessibility is \textbf{not}
required, information and operations can exist without causality.

This corresponds to a state that can exist without possessing causal
structure. The existence without causality is also related to
discussions of ``physics without time'' in quantum gravity {[}4{]}. Such
a state ``exists'' but is ``not cognized,'' and is described without
contradiction under the definitions of §1.

\begin{center}\rule{0.5\linewidth}{0.5pt}\end{center}

\subsection{3. Structural Requirements}\label{structural-requirements}

\subsubsection{3.1 Basic Requirement}\label{basic-requirement}

\textbf{The definitions of information \(I\) and operation \(F^*\) in §1
must be capable of encompassing all information and all operations in
current physics.}

\textbf{Scope:} The ``Theory of Everything'' referred to here targets
the entire domain of physics that is mathematically analyzable. The
humanities and social sciences --- sociology, history, philosophy,
economics, etc. --- are excluded. The structural requirements in this
paper aim to encompass only phenomena describable as physical laws
through mathematical equations, and make no claims regarding
applicability to other domains.

\subsubsection{3.2 Scope of Detailed
Requirements}\label{scope-of-detailed-requirements}

The scope of detailed requirements derived from the basic requirement
(§3.1) is limited to the following:

\begin{itemize}
\tightlist
\item
  Theoretical physics
\item
  Particle physics
\item
  Gravitational field theory
\item
  General relativity
\end{itemize}

Other branches of physics (thermodynamics, statistical mechanics,
condensed matter physics, fluid dynamics, etc.) are excluded from the
detailed requirements of this paper as potential future scope
extensions.

\subsubsection{3.3 Supplementary Definitions Based on
Scope}\label{supplementary-definitions-based-on-scope}

The definitions in §1--§2 are general definitions independent of scope.
The following supplementary definitions are provided to describe the
detailed requirements for the scope of §3.2 (theoretical physics,
particle physics, gravitational field theory, general relativity).

\paragraph{3.3.1 Causality}\label{causality}

The ordering structure that emerged in §2.2 is defined for the present
scope as follows:

\textbf{Causality:} An irreversible ordering relation existing between
the input information \(\mathrm{IN}\) (cause) and the output information
\(\mathrm{OUT}\) (effect) of operation \(F^*\). That is, the direction
\(\mathrm{IN}\) → \(F^*\) → \(\mathrm{OUT}\) is established, and the
reverse direction is not permitted.

The causality defined here is independent of time. It is also
independent of the direction of time (past → future). Time is treated as
a type of information \(I\) as defined in §1.1, and time is not required
for the definition of causality. Causality exists even without the
existence of time.

\paragraph{3.3.2 Conserved Quantities
(Symmetry)}\label{conserved-quantities-symmetry}

\textbf{Conserved quantity:} A quantity of information that remains
invariant before and after operation \(F^*\). That is, a function
\(Q(I)\) of information \(I\) whose value does not change regardless of
how operation \(F^*\) is executed is called a conserved quantity.

By Noether's theorem {[}5{]}, conserved quantities are equivalent to
symmetries. That is, the existence of a conserved quantity \(Q\)
corresponds one-to-one with the invariance of operation \(F^*\) under a
certain transformation (symmetry).

\paragraph{3.3.3 Dimension}\label{dimension}

\textbf{Dimension:} The number of independent directions in which
information \(I\) can vary. That is, the number of independent degrees
of freedom in the state space of information \(I\) is defined as the
dimension.

\begin{itemize}
\tightlist
\item
  \textbf{Null dimension:} A state in which the information of dimension
  itself does not exist (\(d = \mathrm{null}\)). This corresponds to the
  null state of §1.1. Since the concept of dimension itself is not
  defined, nothing can be said about spatial properties.
\item
  \textbf{Zero dimension (\(d = 0\)):} A state in which the dimension
  exists but its value is 0. Information \(I\) has no independent
  direction in which it can vary. However, information itself can exist.
  In zero dimensions, the distinction between ``exists''
  (\(I \neq \mathrm{null}\)) and ``does not exist''
  (\(I = \mathrm{null}\)) can be made. The count or magnitude of
  information is not defined, but whether it exists or not is decidable.
  It should be noted that the definition of information \(I\) in §1.1
  places no restriction on the size or complexity of its content. A
  single piece of information encoding the state of the entire universe
  is also information \(I\). Therefore, even in zero dimensions, the one
  piece of information that ``exists'' can have arbitrarily complex
  content.
\item
  \textbf{\(d\) dimensions (\(d \geq 1\)):} A state in which information
  \(I\) can vary in \(d\) independent directions.
\item
  \textbf{Infinite dimensions:} A state in which there is no upper bound
  on the number of independent directions in which information \(I\) can
  vary.
\end{itemize}

The number of dimensions is determined by the structure of conserved
quantities (symmetries).

\subsubsection{3.4 Detailed Requirements}\label{detailed-requirements}

\paragraph{Requirement 1: Existence of
Causality}\label{requirement-1-existence-of-causality}

\textbf{Requirement:} The theory must possess causal structure. That is,
an ordering must exist between cause and effect.

\textbf{Rationale:} In all of theoretical physics, particle physics,
gravitational field theory, and general relativity, causality is the
most fundamental premise. No physical theory exists without causality.

\textbf{Sufficiency:} By §2, causal structure inevitably emerges at the
point when cognitive accessibility is required. Therefore, this
requirement is already satisfied by the definitions of §1 and the
requirement of cognitive accessibility in §2.

\paragraph{Requirement 2: Conserved Quantities Must Be
Definable}\label{requirement-2-conserved-quantities-must-be-definable}

\textbf{Requirement:} The theory must be able to define quantities
(conserved quantities) that remain invariant before and after operation
\(F^*\).

\textbf{Rationale:} In every branch of physics, conserved quantities are
the most fundamental structures. Physical theories cannot exist without
conserved quantities such as energy conservation, momentum conservation,
angular momentum conservation, and charge conservation. By Noether's
theorem {[}5{]}, conserved quantities are equivalent to symmetries.
Symmetries generate conserved quantities, and the structure of conserved
quantities determines the dimensions of space. Therefore, the
definability of conserved quantities precedes the existence of
dimensions.

\textbf{Sufficiency:} Operation \(F^*\) of §1 has the rule that when
\(\mathrm{IN} = I\), it does nothing (\(\mathrm{OUT}\) remains
unchanged). This means that when the input matches the internal state,
the output is conserved. Conserved quantities can be defined as
combinations of information satisfying this invariance condition.

\paragraph{Requirement 3: Dimensions Must Be
Definable}\label{requirement-3-dimensions-must-be-definable}

\textbf{Requirement:} The theory must possess the concept of dimension
as defined in §3.3.3. That is, it must be possible to define the
dimension as the number of independent degrees of freedom in the state
space of information \(I\).

\textbf{Scope dependency:} This requirement depends on the scope of
§3.2. The general definitions of §1--§2 do not presuppose the existence
of dimensions, but within the scope of this paper (theoretical physics,
particle physics, gravitational field theory, general relativity), all
fields are formulated upon dimensional structure. Theories that do not
require dimensions are outside the scope of this paper. Therefore, the
definability of dimensions is a mandatory requirement within the present
scope.

\textbf{Sufficiency:} By §3.3.3, dimension is definable as the number of
independent directions in which information \(I\) can vary. Null
dimension, zero dimension, \(d\) dimensions, and infinite dimensions are
each clearly distinguished.

\textbf{Extended requirements:} For a Theory of Everything, the
following requirements --- which are not mandatory within the current
scope --- are defined as necessary:

\begin{itemize}
\tightlist
\item
  \textbf{Infinite extensibility of dimensions:} Current theories
  presuppose a specific number of dimensions (such as 3+1). However, a
  Theory of Everything is required to be non-degenerate and regular for
  any number of dimensions from zero to infinity.
\item
  \textbf{Symmetry of dimensions:} Current theories assign specific
  meanings to particular axes (such as the distinction between the time
  axis and spatial axes). However, a Theory of Everything requires that
  all dimensions are equivalent (that dimensions possess symmetry).
  Assigning a special role to a particular axis must be derived as a
  result of spontaneous symmetry breaking and must not be included in
  the definition.
\item
  \textbf{Indifference to continuous/discrete:} Current theories
  presuppose that dimensions are continuous. However, a Theory of
  Everything is required to be non-degenerate and regular whether
  dimensions take continuous or discrete values.
\end{itemize}

\paragraph{Requirement 4: Reversibility of
Causality}\label{requirement-4-reversibility-of-causality}

\textbf{Requirement:} The theory must be non-degenerate and regular even
when the direction of causality is reversed. That is, even when the
direction \(\mathrm{IN}\) (cause) → \(F^*\) → \(\mathrm{OUT}\) (effect)
is reinterpreted as \(\mathrm{OUT}\) (cause) → \(F^*\) → \(\mathrm{IN}\)
(effect), the structure of the theory must be preserved.

\textbf{Scope dependency:} This requirement depends on the scope of
§3.2. In all four fields of the present scope (theoretical physics,
particle physics, gravitational field theory, general relativity), the
fundamental equations are all symmetric under reversal of causality.
Theories that are asymmetric under reversal of causality (such as
thermodynamics) are outside the present scope.

\textbf{Distinction between reversal of causality and time reversal:}
What is required here is strictly the \textbf{reversal of causality},
which is not synonymous with time reversal. As defined in §3.3.1,
causality is defined independently of time. Whether time reversal is
required when causality is reversed depends on the specific construction
of operation \(F^*\) in individual theories, and is not included in the
structural requirements of this paper.

\textbf{Sufficiency:} The definition of operation \(F^*\) in §1.2 does
not prohibit the exchange of the roles of input and output. The reversal
of causality corresponds to reversing the direction while preserving the
cognitive accessibility (existence of ordering) of §2.1.

\paragraph{Requirement 5: Inverse Operations Must Be
Definable}\label{requirement-5-inverse-operations-must-be-definable}

\textbf{Requirement:} The theory must be able to define an inverse
operation that undoes the result of operation \(F^*\).

\textbf{Scope dependency:} This requirement depends on the scope of
§3.2. In theoretical physics, the reversibility of unitary
transformations; in particle physics, CPT symmetry; and in gravitational
field theory and general relativity, the reversibility of field
equations --- each requires the existence of inverse operations.

\textbf{Sufficiency:} Operation \(F^*\) of §1.2 is defined as a mapping
from information to information and does not prohibit inverse
operations. An inverse operation can be defined as a separate operation
from \(F^*\).

\paragraph{Requirement 6: Composition of Operations Must Be
Possible}\label{requirement-6-composition-of-operations-must-be-possible}

\textbf{Requirement:} The theory must be able to chain and compose
multiple operations \(F^*\). That is, it must be possible to use the
output of one operation as the input to another operation.

\textbf{Scope dependency:} This requirement depends on the scope of
§3.2. In theoretical physics, the group structure of symmetries; in
particle physics, field interactions and perturbation expansions; in
gravitational field theory, nonlinear field equations; and in general
relativity, the composition of coordinate transformations --- each
requires the composition of operations.

\textbf{Sufficiency:} Operation \(F^*\) of §1.2 requires that input
information and output information belong to the same type. Therefore,
the output of one operation can be passed as input to another operation
of the same type, and the composition of operations is guaranteed by
definition.

\paragraph{Requirement 7: Finiteness of Values Must Be
Guaranteeable}\label{requirement-7-finiteness-of-values-must-be-guaranteeable}

\textbf{Requirement:} When the values of information \(I\) can become
infinite, the theory must have a method to map them to a finite range.
That is, it must be possible to structurally avoid the divergence of the
values of information \(I\) through repeated application of operation
\(F^*\).

\textbf{Scope dependency:} This requirement depends on the scope of
§3.2. In theoretical physics and particle physics, ultraviolet and
infrared divergences in quantum field theory are fundamental problems
addressed by renormalization theory {[}6{]}. In gravitational field
theory and general relativity, divergences at singularities (such as the
center of black holes) remain unsolved problems {[}7{]}. All of these
stem from the fact that finiteness of values is not structurally
guaranteed. A Theory of Everything requires the existence of a mapping
that guarantees finiteness of values at the structural level, rather
than relying on individual techniques (such as renormalization).

\textbf{Sufficiency:} The definition of information \(I\) in §1.1 does
not restrict the range of values. Therefore, there is room to construct
a mapping to a finite range. However, the specific method of mapping
(embedding into phase space, compactification, etc.) is outside the
scope of the structural requirements of this paper and is left to the
construction of individual theories.

\textbf{Regularity of the inverse mapping:} Furthermore, the inverse
mapping of this finitization mapping must also be regular. That is, the
original information must be recoverable from the information mapped to
a finite range. This is already required by Requirement 5 (existence of
inverse operations) and the definition of operation \(F^*\) in §1.2
(input information and output information belong to the same type), but
is restated here in the context of finiteness of values. If a
finitization mapping irreversibly loses information, it does not satisfy
the requirements of operation \(F^*\).

\paragraph{Requirement 8: Dimensional
Unification}\label{requirement-8-dimensional-unification}

\textbf{Requirement:} The dimensions of all physical quantities
appearing in the theory must be expressible as powers \(L^n\) of a
single fundamental dimension \(L\) (length).

\textbf{Rationale:} In current physics, four independent dimensions are
used: mass \(M\), length \(L\), time \(T\), and charge \(Q\). However,
upon adopting natural units (\(c = 1\), \(\hbar = 1\)):

\[[T] = [L], \quad [M] = [L^{-1}], \quad [E] = [L^{-1}]\]

The dimensions of major physical quantities are expressed in terms of
\(L^n\) as follows:

{\def\LTcaptype{none} % do not increment counter
\begin{longtable}[]{@{}
  >{\raggedright\arraybackslash}p{(\linewidth - 4\tabcolsep) * \real{0.3333}}
  >{\raggedright\arraybackslash}p{(\linewidth - 4\tabcolsep) * \real{0.3333}}
  >{\raggedright\arraybackslash}p{(\linewidth - 4\tabcolsep) * \real{0.3333}}@{}}
\toprule\noalign{}
\begin{minipage}[b]{\linewidth}\raggedright
Physical quantity
\end{minipage} & \begin{minipage}[b]{\linewidth}\raggedright
SI dimension
\end{minipage} & \begin{minipage}[b]{\linewidth}\raggedright
\(L^n\) expression
\end{minipage} \\
\midrule\noalign{}
\endhead
\bottomrule\noalign{}
\endlastfoot
Length \(L\) & \(L\) & \(L^1\) \\
Time \(T\) & \(T\) & \(L^1\) (\(c = 1\)) \\
Mass \(M\) & \(M\) & \(L^{-1}\) (\(\hbar = 1\)) \\
Energy \(E\) & \(ML^2T^{-2}\) & \(L^{-1}\) \\
Velocity \(v\) & \(LT^{-1}\) & \(L^0\) (dimensionless) \\
Gravitational constant \(G\) & \(M^{-1}L^3T^{-2}\) & \(L^2\) \\
Planck constant \(\hbar\) & \(ML^2T^{-1}\) & \(L^0\) (by definition,
\(\hbar = 1\)) \\
Speed of light \(c\) & \(LT^{-1}\) & \(L^0\) (by definition,
\(c = 1\)) \\
\end{longtable}
}

\textbf{Consequence for the charge dimension:} Under this requirement,
charge \(Q\) must also be expressible as \(L^n\). In Gaussian units,
\([Q] = [M^{1/2}L^{3/2}T^{-1}]\), and under \(c = \hbar = 1\), this
yields \([Q] = [L^0]\) (dimensionless). That is, the elementary charge
\(e\) becomes a dimensionless pure numerical constant, and the
fine-structure constant \(\alpha = e^2/(4\pi)\) also becomes a pure
geometric constant.

\textbf{Verification:} The major equations of gravitational field
theory, electromagnetism, quantum mechanics, and thermodynamics all
satisfy this requirement.

\textbf{Sufficiency:} Since the definition of information \(I\) in §1.1
does not constrain the structure of dimensions, a construction that
unifies all dimensions into \(L^n\) is permitted.

\paragraph{Requirement 9: Scale
Symmetry}\label{requirement-9-scale-symmetry}

\textbf{Requirement:} All equations appearing in the theory must be
scale-homogeneous under the scale transformation \(L \to \lambda L\)
(\(\lambda > 0\)). That is, when each physical quantity has dimension
\(L^n\), it transforms as \(L^n \to \lambda^n L^n\) under
\(L \to \lambda L\), and both sides of any equation must carry the same
power \(\lambda^k\).

\textbf{Rationale:} For a Theory of Everything to depend on a specific
scale (a specific length) would mean that another theory is needed to
explain the origin of that scale. A truly fundamental theory must be
scale-homogeneous (power-law). Logarithmic scale dependence (e.g., the
logarithmic running of the coupling constant \(\alpha_s(\mu)\) in
quantum chromodynamics) does not satisfy this requirement.

\textbf{Formulation:} By Requirement 8, all physical quantities have
dimension \(L^n\). Under the scale transformation \(L \to \lambda L\):

\[f(\lambda L) = \lambda^k f(L)\]

is required to hold. This means that \(f\) is a homogeneous function of
degree \(k\) in \(L\).

\textbf{Verification results:} This requirement was verified for the
following major equations:

\begin{itemize}
\tightlist
\item
  \textbf{Classical mechanics:} Newton's equation of motion \(F = ma\),
  gravitational force \(F = GMm/r^2\), Kepler's third law
  \(T^2 \propto a^3\) --- all scale-homogeneous ✓
\item
  \textbf{Gravitational field theory:} Einstein equation
  \(G_{\mu\nu} = 8\pi G T_{\mu\nu}\), Schwarzschild solution
  \(r_s = 2GM/c^2\), geodesic equation --- all scale-homogeneous ✓
\item
  \textbf{Electromagnetism:} Maxwell's equations, Coulomb force
  \(F = e^2/(4\pi r^2)\), Lorentz force --- all scale-homogeneous ✓
\item
  \textbf{Quantum mechanics:} Schrödinger equation, Dirac equation,
  Klein--Gordon equation --- all scale-homogeneous ✓
\item
  \textbf{Thermodynamics:} Stefan--Boltzmann law \(P \propto T^4\),
  blackbody radiation spectrum --- all scale-homogeneous ✓
\end{itemize}

\textbf{Exception:} The coupling constant \(\alpha_s(\mu)\) of quantum
chromodynamics (QCD) exhibits logarithmic running and does not satisfy
this requirement. This suggests that QCD is an effective theory
incorporating renormalization group effects and requires derivation from
a more fundamental theory.

\textbf{Sufficiency:} Since the definitions of §1 do not prescribe any
specific behavior under scale transformations, the construction of a
scale-homogeneous theory is permitted.

\begin{center}\rule{0.5\linewidth}{0.5pt}\end{center}

\subsection{4. Remarks: Examples of Constructions That Facilitate
Satisfaction of Structural
Requirements}\label{remarks-examples-of-constructions-that-facilitate-satisfaction-of-structural-requirements}

This chapter does not constitute mandatory requirements. It lists
techniques believed to facilitate the satisfaction of the structural
requirements of §3 when constructing a theory.

\subsubsection{4.1 Central Projection}\label{central-projection}

Using central projection as a mapping from coordinate space to phase
space facilitates the satisfaction of Requirement 7 (finiteness of
values). Central projection is a method that maps a coordinate space
including infinity to a finite phase space, and the inverse mapping is
also regular.

\subsubsection{4.2 Treating Coordinate Information as Phase
Information}\label{treating-coordinate-information-as-phase-information}

When coordinate information is treated as phase information (angles),
the domain of values becomes bounded (e.g., \(0\) to \(2\pi\)), making
it easier to structurally guarantee Requirement 7 (finiteness of
values).

\begin{center}\rule{0.5\linewidth}{0.5pt}\end{center}

\subsection{References}\label{references}

{[}1{]} C. E. Shannon, ``A Mathematical Theory of Communication,''
\emph{Bell System Technical Journal}, vol.~27, pp.~379--423, 623--656,
1948. --- Mathematical foundations of the definition of information.

{[}2{]} A. M. Turing, ``On Computable Numbers, with an Application to
the Entscheidungsproblem,'' \emph{Proceedings of the London Mathematical
Society}, ser. 2, vol.~42, pp.~230--265, 1937. --- Formal definition of
computability and operations.

{[}3{]} C. J. Isham, ``Canonical Quantum Gravity and the Problem of
Time,'' in \emph{Integrable Systems, Quantum Groups, and Quantum Field
Theories}, NATO ASI Series, vol.~409, pp.~157--287, 1993. --- The frozen
time problem in quantum gravity; the absence of a global time parameter.

{[}4{]} C. Rovelli, ``Forget Time,'' \emph{Foundations of Physics},
vol.~41, pp.~1475--1490, 2011. arXiv:0903.3832. --- The possibility of
``timeless'' formulations in fundamental physics.

{[}5{]} E. Noether, ``Invariante Variationsprobleme,'' \emph{Nachrichten
von der Gesellschaft der Wissenschaften zu Göttingen}, pp.~235--257,
1918. --- Equivalence of symmetries and conserved quantities (Noether's
theorem).

{[}6{]} S. Weinberg, ``Living with Infinities,'' arXiv:0903.0568, 2009.
--- Overview of divergence problems and renormalization in quantum field
theory.

{[}7{]} R. Penrose, ``Gravitational Collapse and Space-Time
Singularities,'' \emph{Physical Review Letters}, vol.~14, pp.~57--59,
1965. --- Singularity theorems in general relativity.

\end{document}
